\begin{theorem}[Hankel 指纹的有限模单值化与最短锁定阈值]\label{thm:conclusion-hankel-finite-mod-singlevaluedization-shortest-locking-threshold}
\leanverified{paper\_conclusion\_hankel\_finite\_mod\_singlevaluedization\_shortest\_locking\_threshold}
设整数序列 \((a_n)_{n\ge 0}\subset \ZZ\) 的 Hankel 秩为 \(d\)，并记连续 \(d\times d\) Hankel 主块
\[
H_k:=\bigl(a_{k+i+j}\bigr)_{0\le i,j\le d-1}.
\]
若存在某个可逆见证块
\[
H_{k_0}\in M_d(\ZZ),
\qquad
\det(H_{k_0})\neq 0,
\]
则对任意满足
\[
\gcd(\det(H_{k_0}),M)=1
\]
的模数 \(M\ge 2\)，矩阵
\[
\bar A:=\bar H_{k_0}^{-1}\bar H_{k_0+1}\in GL_d(\ZZ/M\ZZ)
\]
良定义，并且全部后续 Hankel 轨道满足严格闭式
\[
\bar H_{k_0+r}=\bar H_{k_0}\bar A^r
\qquad (r\ge 0).
\]
因此模 \(M\) 的 Hankel 未来轨道是纯周期序列，其周期整除 \(\bar A\) 在 \(GL_d(\ZZ/M\ZZ)\) 中的阶；特别地，良模数层面的全部未来仅由一对相邻块
\[
(\bar H_{k_0},\bar H_{k_0+1})
\]
唯一决定。

另一方面，该有限锁定阈值在确定性意义下不可再压缩：一般位置下，长度 \(2d-1\) 的连续窗口不足以唯一确定全序列，而长度 \(2d\) 的窗口恰是最短锁定阈值。
\end{theorem}

\begin{proof}
正文关于最小 Hankel realization 的传播公式已经证明：当 Hankel 秩为 \(d\) 且 \(H_{k_0}\) 可逆时，存在唯一偏移传递矩阵
\[
A:=H_{k_0}^{-1}H_{k_0+1}
\]
使得
\[
H_{k_0+r}=H_{k_0}A^r
\qquad (r\ge 0).
\]
在模 \(M\) 上约化后，由 \(\gcd(\det(H_{k_0}),M)=1\) 可知 \(\bar H_{k_0}\) 可逆，于是
\[
\bar A=\bar H_{k_0}^{-1}\bar H_{k_0+1}
\]
良定义，并保留同一传播公式
\[
\bar H_{k_0+r}=\bar H_{k_0}\bar A^r.
\]
由于 \(GL_d(\ZZ/M\ZZ)\) 为有限群，\(\bar A\) 具有有限阶 \(\operatorname{ord}(\bar A)\)，故
\[
\bar H_{k_0+r+\operatorname{ord}(\bar A)}
=
\bar H_{k_0}\bar A^{r+\operatorname{ord}(\bar A)}
=
\bar H_{k_0}\bar A^r
=
\bar H_{k_0+r},
\]
于是未来轨道纯周期，且周期整除 \(\operatorname{ord}(\bar A)\)。

至于阈值尖锐性，正文已经证明：长度 \(2d\) 的连续窗口恰可唯一锁定最小 Hankel realization，而长度 \(2d-1\) 的连续窗口在一般位置下仍保留一维自由度。故最短锁定阈值正是 \(2d\)。证毕。
\end{proof}

\begin{corollary}[Hankel 模退化的有限坏素数局域化]\label{cor:conclusion-hankel-mod-degeneration-finite-bad-prime-localization}
沿用定理 \ref{thm:conclusion-hankel-finite-mod-singlevaluedization-shortest-locking-threshold} 的记号，全部后续 Hankel 行列式的素因子集合满足
\[
\bigcup_{r\ge 0}\{p:\ p\mid \det(H_{k_0+r})\}
\subseteq
\{p:\ p\mid \det(H_{k_0})\det(H_{k_0+1})\}.
\]
因此，模退化只能发生在一个有限坏素数集上；对所有其余素数，Hankel 正规形与未来轨道都处于良模数单值化状态。
\leanverified{paper\_conclusion\_hankel\_mod\_degeneration\_finite\_bad\_prime\_localization}
\end{corollary}

\begin{proof}
由定理 \ref{thm:conclusion-hankel-finite-mod-singlevaluedization-shortest-locking-threshold} 的传播公式，
\[
H_{k_0+r}=H_{k_0}A^r,
\qquad
A=H_{k_0}^{-1}H_{k_0+1}.
\]
取行列式即得
\[
\det(H_{k_0+r})=\det(H_{k_0})\det(A)^r.
\]
而
\[
\det(A)=\frac{\det(H_{k_0+1})}{\det(H_{k_0})}.
\]
故任一未来行列式 \(\det(H_{k_0+r})\) 的素因子都只能来自 \(\det(H_{k_0})\) 与 \(\det(H_{k_0+1})\) 的素因子。于是坏素数集有限，并完全局域于初始相邻两块之上。证毕。
\end{proof}
